\documentclass{article}
\usepackage[utf8]{inputenc}
\usepackage[italian]{babel}
\usepackage{geometry}
\usepackage{graphicx}
\usepackage{booktabs}
\usepackage{longtable}
\usepackage{hyperref}
\usepackage{xcolor}
\usepackage{enumitem}
\usepackage{fancyhdr}
\usepackage{listings}
\usepackage{verbatim}

\geometry{margin=2.5cm}

\hypersetup{
    colorlinks=true,
    linkcolor=blue,
    filecolor=magenta,      
    urlcolor=cyan,
    pdftitle={Documentazione Architettura KARIS},
    pdfauthor={Tommaso Caputi, Di Pasquale Giulio, Gadaleta Giovanni}
}

\pagestyle{fancy}
\fancyhf{}
\fancyhead[L]{Documentazione Architettura - KARIS}
\fancyhead[R]{\thepage}
\fancyfoot[C]{}

\title{Documentazione Architettura e Tech Stack - KARIS}
\author{Tommaso Caputi, Di Pasquale Giulio, Gadaleta Giovanni}

% Configurazione per codice
\lstset{
    basicstyle=\ttfamily\small,
    breaklines=true,
    frame=single,
    numbers=left,
    numberstyle=\tiny,
    showstringspaces=false,
    extendedchars=true,
    escapeinside={\%*}{*)},
    inputencoding=utf8
}

\begin{document}

\maketitle
\newpage

\tableofcontents
\newpage

\section{Introduzione}

\textbf{KARIS} è un'applicazione web full-stack progettata per la gestione digitale delle risorse e della distribuzione di beni nelle strutture Caritas parrocchiali. Il sistema permette di gestire inventari, assegnazioni, beneficiari e facilita l'interscambio cooperativo tra diverse parrocchie.

Questa documentazione descrive l'architettura tecnica del sistema, le tecnologie utilizzate, come i componenti interagiscono tra loro e presenta un caso d'uso completo con il relativo workflow.

\section{Tech Stack}

\subsection{Frontend}

\begin{longtable}{p{4cm}p{3cm}p{8cm}}
\toprule
\textbf{Tecnologia} & \textbf{Versione} & \textbf{Utilizzo} \\
\midrule
\endfirsthead
\toprule
\textbf{Tecnologia} & \textbf{Versione} & \textbf{Utilizzo} \\
\midrule
\endhead
\bottomrule
\endfoot
\bottomrule
\endlastfoot
\textbf{Next.js} & 16.0.10 & Framework React per SSR/SSG e routing \\
\textbf{React} & 19.2.1 & Libreria UI per componenti interattivi \\
\textbf{TypeScript} & 5.x & Linguaggio per type safety \\
\textbf{Tailwind CSS} & 4.x & Framework CSS utility-first \\
\textbf{Radix UI} & Latest & Componenti UI accessibili e headless \\
\textbf{Lucide React} & 0.561.0 & Libreria di icone \\
\textbf{Sonner} & 2.0.7 & Sistema di notifiche toast \\
\end{longtable}

\subsection{Backend}

\begin{longtable}{p{4cm}p{3cm}p{8cm}}
\toprule
\textbf{Tecnologia} & \textbf{Versione} & \textbf{Utilizzo} \\
\midrule
\endfirsthead
\toprule
\textbf{Tecnologia} & \textbf{Versione} & \textbf{Utilizzo} \\
\midrule
\endhead
\bottomrule
\endfoot
\bottomrule
\endlastfoot
\textbf{Next.js API Routes} & 16.0.10 & Endpoint REST per logica server-side \\
\textbf{Supabase JS Client} & 2.49.1 & Client per database PostgreSQL e autenticazione \\
\end{longtable}

\subsection{Database e Infrastruttura}

\begin{longtable}{p{4cm}p{11cm}}
\toprule
\textbf{Tecnologia} & \textbf{Utilizzo} \\
\midrule
\endfirsthead
\toprule
\textbf{Tecnologia} & \textbf{Utilizzo} \\
\midrule
\endhead
\bottomrule
\endfoot
\bottomrule
\endlastfoot
\textbf{PostgreSQL} & Database relazionale (gestito via Supabase) \\
\textbf{Supabase} & BaaS (Backend as a Service) per database e auth \\
\textbf{Supabase Auth} & Sistema di autenticazione e gestione utenti \\
\end{longtable}

\subsection{Strumenti di Sviluppo}

\begin{longtable}{p{4cm}p{11cm}}
\toprule
\textbf{Tecnologia} & \textbf{Utilizzo} \\
\midrule
\endfirsthead
\toprule
\textbf{Tecnologia} & \textbf{Utilizzo} \\
\midrule
\endhead
\bottomrule
\endfoot
\bottomrule
\endlastfoot
\textbf{ESLint} & Linter per qualità del codice \\
\textbf{PostCSS} & Processore CSS \\
\textbf{TypeScript Compiler} & Compilazione e type checking \\
\end{longtable}

\section{Architettura del Sistema}

\subsection{Architettura Generale}

KARIS segue un'architettura \textbf{full-stack monolitica} basata su Next.js, che integra frontend e backend in un'unica applicazione. Il sistema utilizza il pattern \textbf{App Router} di Next.js 13+ per il routing e la gestione delle pagine.

Il sistema è organizzato in più livelli:

\begin{enumerate}[leftmargin=*]
\item \textbf{Browser Client}: Componenti React renderizzati nel browser
\item \textbf{Next.js App Router}: Server Components per routing e rendering
\item \textbf{Next.js API Routes}: Endpoint REST per logica backend
\item \textbf{Supabase Client}: Interfaccia per database e autenticazione
\item \textbf{Supabase Cloud}: Database PostgreSQL e servizi di autenticazione
\end{enumerate}

\subsection{Pattern Architetturali}

\begin{enumerate}[leftmargin=*]
\item \textbf{Server-Side Rendering (SSR)}: Le pagine vengono renderizzate sul server per migliorare SEO e performance iniziale
\item \textbf{Client-Side Rendering (CSR)}: Componenti interattivi renderizzati lato client
\item \textbf{API Routes}: Endpoint REST per operazioni CRUD e logica di business
\item \textbf{Component-Based Architecture}: UI costruita con componenti React riutilizzabili
\item \textbf{Separation of Concerns}: Separazione tra logica di presentazione, business e data access
\end{enumerate}

\section{Struttura del Progetto}

La struttura del progetto è organizzata come segue:

\begin{verbatim}
karis/
+-- db/                          # Script SQL per database
|   +-- schema.sql              # Schema completo del database
|   +-- seed.sql                # Dati di esempio
|
+-- documentazione/              # Documentazione del progetto
|   +-- DOCUMENTAZIONE_DB.md
|   +-- DOCUMENTAZIONE_ARCHITETTURA.md
|   +-- er_schema.png
|
+-- karis/                       # Applicazione Next.js
    +-- src/
    |   +-- app/                 # App Router (Next.js 13+)
    |   |   +-- api/             # API Routes (Backend)
    |   |   |   +-- assegnazioni/
    |   |   |   +-- beneficiario/
    |   |   |   +-- beni/
    |   |   |   +-- dashboard/
    |   |   |   +-- famiglia/
    |   |   |   +-- inviti/
    |   |   |   +-- richieste/
    |   |   |   +-- user/
    |   |   |   +-- volontari/
    |   |   |
    |   |   +-- beneficiario/    # Pagine per gestione beneficiari
    |   |   +-- beni/            # Pagine per gestione beni
    |   |   +-- dashboard/       # Dashboard principale
    |   |   +-- famiglia/        # Pagine per gestione famiglie
    |   |   +-- invito/          # Pagina per accettazione inviti
    |   |   +-- login/           # Pagina di login
    |   |   +-- richieste/       # Pagina per richieste tra parrocchie
    |   |   +-- volontari/       # Pagina per gestione volontari
    |   |   +-- layout.tsx       # Layout root dell'applicazione
    |   |   +-- page.tsx         # Homepage (landing page)
    |   |   +-- globals.css      # Stili globali
    |   |
    |   +-- components/          # Componenti React riutilizzabili
    |   |   +-- ui/              # Componenti UI base (shadcn/ui)
    |   |   |   +-- button.tsx
    |   |   |   +-- dialog.tsx
    |   |   |   +-- input.tsx
    |   |   |   +-- ...
    |   |   +-- DashboardLayout.tsx    # Layout condiviso per dashboard
    |   |   +-- LandingPage/          # Componenti per landing page
    |   |   +-- NotFound.tsx
    |   |
    |   +-- lib/                 # Utilities e helper
    |       +-- supabaseClient.ts     # Client Supabase configurato
    |       +-- apiHelpers.ts         # Helper per API routes
    |       +-- utils.ts              # Utility generiche
    |
    +-- public/                  # File statici (immagini, favicon, etc.)
    +-- package.json
    +-- tsconfig.json
    +-- next.config.ts
    +-- tailwind.config.js
\end{verbatim}

\section{Componenti e Interazioni}

\subsection{Livello 1: Componenti UI (Frontend)}

\subsubsection{DashboardLayout}
\textbf{Posizione}: \texttt{src/components/DashboardLayout.tsx}

\textbf{Responsabilità}:
\begin{itemize}
\item Layout condiviso per tutte le pagine autenticate
\item Sidebar di navigazione
\item Header con informazioni utente
\item Gestione autenticazione e logout
\end{itemize}

\textbf{Interazioni}:
\begin{itemize}
\item Utilizza \texttt{supabase.auth.getUser()} per verificare autenticazione
\item Chiama \texttt{/api/user} per recuperare dati utente
\item Gestisce navigazione tra sezioni
\end{itemize}

\subsubsection{Componenti di Pagina}
Ogni pagina (es. \texttt{beni/page.tsx}, \texttt{beneficiario/page.tsx}) è un componente React che:
\begin{itemize}
\item Utilizza \texttt{DashboardLayout} come wrapper
\item Effettua chiamate API per recuperare dati
\item Gestisce stato locale con React hooks (\texttt{useState}, \texttt{useEffect})
\item Mostra feedback all'utente tramite toast notifications (Sonner)
\end{itemize}

\subsection{Livello 2: API Routes (Backend)}

\subsubsection{Struttura API Route}
Ogni API route segue questo pattern:

\begin{lstlisting}
// src/app/api/[resource]/route.ts
export async function GET(request: Request) {
    // 1. Validazione parametri
    // 2. Autenticazione utente
    // 3. Recupero dati da database
    // 4. Formattazione risposta
    // 5. Return JSON
}

export async function POST(request: Request) {
    // 1. Validazione body
    // 2. Autenticazione utente
    // 3. Validazione business logic
    // 4. Inserimento/aggiornamento database
    // 5. Return risultato
}
\end{lstlisting}

\subsubsection{Helper Functions}
\textbf{Posizione}: \texttt{src/lib/apiHelpers.ts}

Funzioni riutilizzabili per API routes:
\begin{itemize}
\item \texttt{getUserParrocchia(userId)}: Recupera utente e parrocchia associata
\item \texttt{validateUserId()}: Valida presenza userId nei parametri
\item \texttt{normalizeSupabaseRelation()}: Normalizza relazioni Supabase
\end{itemize}

\subsection{Livello 3: Database Layer}

\subsubsection{Supabase Client}
\textbf{Posizione}: \texttt{src/lib/supabaseClient.ts}

Configurazione del client Supabase:
\begin{itemize}
\item Inizializzazione con URL e API key da variabili d'ambiente
\item Gestione errori quando Supabase non è configurato
\item Utilizzato sia lato client che server
\end{itemize}

\subsubsection{Query Pattern}
Le query seguono questo pattern:

\begin{lstlisting}
const { data, error } = await supabase
    .from("tabella")
    .select("campo1, campo2, relazione:campo_id (id, nome)")
    .eq("campo", valore)
    .order("created_at", { ascending: false });
\end{lstlisting}

\subsection{Flusso di Dati Completo}

Il flusso di dati segue questa sequenza:

\begin{enumerate}[leftmargin=*]
\item \textbf{User Action}: L'utente interagisce con il componente (click, submit)
\item \textbf{React Component}: Gestisce lo stato e prepara la richiesta
\item \textbf{API Call}: Chiamata \texttt{fetch('/api/...')} al server
\item \textbf{API Route}: Validazione, autenticazione e business logic
\item \textbf{Supabase Client}: Query builder per database
\item \textbf{Database}: Operazioni su PostgreSQL
\item \textbf{Response}: Ritorno dati formattati in JSON
\item \textbf{UI Update}: Aggiornamento stato e re-render
\end{enumerate}

\section{Flusso di Autenticazione}

\subsection{Login}

Il processo di login segue questi passaggi:

\begin{enumerate}[leftmargin=*]
\item L'utente inserisce email e password nel form di login
\item Il componente \texttt{/login/page.tsx} gestisce lo stato del form
\item Chiamata a \texttt{supabase.auth.signInWithPassword()}
\item Supabase Auth verifica le credenziali e genera un session token
\item Il client salva la session in localStorage
\item Redirect automatico alla dashboard
\end{enumerate}

\subsection{Verifica Autenticazione}

Ogni pagina protetta verifica l'autenticazione:

\begin{lstlisting}
// Pattern comune in tutte le pagine
useEffect(() => {
    const checkAuth = async () => {
        const { data: { user }, error } = await supabase.auth.getUser();
        if (error || !user) {
            router.push("/login");
        }
    };
    checkAuth();
}, []);
\end{lstlisting}

\subsection{Recupero Dati Utente}

Dopo l'autenticazione, il sistema recupera i dati dell'utente dal database:

\begin{lstlisting}
// Chiamata API
const res = await fetch(`/api/user?userId=${user.id}`);
const userData = await res.json();
// userData contiene: nome, cognome, parrocchia, ruolo
\end{lstlisting}

\subsection{Autorizzazione}

Le API routes verificano che l'utente abbia accesso alle risorse della propria parrocchia:

\begin{lstlisting}
// Pattern in API routes
const userResult = await getUserParrocchia(userId);
const parrocchiaId = userResult.data.parrocchia_id;

// Filtra query per parrocchia
query = query.eq("parrocchia_id", parrocchiaId);
\end{lstlisting}

\section{Caso d'Uso: Assegnazione di un Bene}

\subsection{Scenario}

Un volontario della Parrocchia di San Paolo vuole assegnare 5 scatolette di tonno a un beneficiario registrato nel sistema.

\subsection{Workflow Completo}

\subsubsection{Fase 1: Accesso alla Funzionalità}

\begin{enumerate}[leftmargin=*]
\item Il volontario naviga a \texttt{/beni}
\item La pagina mostra la lista dei beni disponibili con pulsante "Assegna" per ogni bene
\item Click su "Assegna" per il tonno
\item Navigazione a \texttt{/beni/assegna?beneId=...}
\item Caricamento dati iniziali (beni, beneficiari, famiglie)
\end{enumerate}

\textbf{Codice rilevante}:
\begin{lstlisting}
// /beni/assegna/page.tsx
useEffect(() => {
    const loadData = async () => {
        // 1. Verifica autenticazione
        const { data: { user } } = await supabase.auth.getUser();
        
        // 2. Carica beni disponibili
        const beniRes = await fetch(`/api/beni?userId=${user.id}`);
        const beni = await beniRes.json();
        
        // 3. Carica beneficiari
        const beneficiariRes = await fetch(`/api/beneficiario?userId=${user.id}`);
        const beneficiari = await beneficiariRes.json();
        
        // 4. Carica famiglie
        const famiglieRes = await fetch(`/api/famiglia?userId=${user.id}&all=true`);
        const famiglie = await famiglieRes.json();
    };
    loadData();
}, []);
\end{lstlisting}

\subsubsection{Fase 2: Selezione e Compilazione Form}

Il form di assegnazione richiede:
\begin{enumerate}[leftmargin=*]
\item Selezione Bene: Tonno (già pre-selezionato)
\item Tipo: Beneficiario o Famiglia (utente seleziona)
\item Beneficiario: Mario Rossi (utente seleziona)
\item Quantità: 5 (utente inserisce)
\item Note: (opzionale)
\end{enumerate}

\textbf{Validazione Client-Side}:
\begin{lstlisting}
// Verifica quantita disponibile
const beneSelezionato = beni.find(b => b.id === formData.beneId);
if (quantitaNum > beneSelezionato.quantity) {
    toast.error(`Quantita non disponibile. Disponibile: ${beneSelezionato.quantity}`);
    return;
}
\end{lstlisting}

\subsubsection{Fase 3: Invio Richiesta}

Quando l'utente clicca "Assegna Bene":
\begin{enumerate}[leftmargin=*]
\item Validazione del form
\item Preparazione del payload
\item Invio POST a \texttt{/api/assegnazioni}
\end{enumerate}

\subsubsection{Fase 4: Elaborazione Backend}

L'API route \texttt{/api/assegnazioni/route.ts} (POST) esegue:

\begin{enumerate}[leftmargin=*]
\item \textbf{Validazione Parametri}:
   \begin{itemize}
   \item \texttt{risorsa\_id} presente?
   \item \texttt{quantita} $>$ 0?
   \item \texttt{beneficiario\_id} o \texttt{famiglia\_id} presente?
   \end{itemize}

\item \textbf{Autenticazione}:
   \begin{itemize}
   \item \texttt{getUserParrocchia(userId)}
   \item Recupera \texttt{parrocchia\_id} utente
   \end{itemize}

\item \textbf{Verifica Risorsa}:
   \begin{itemize}
   \item Risorsa esiste?
   \item Risorsa appartiene alla parrocchia?
   \end{itemize}

\item \textbf{Verifica Beneficiario}:
   \begin{itemize}
   \item Beneficiario esiste?
   \item Beneficiario appartiene alla parrocchia?
   \end{itemize}

\item \textbf{Controllo Disponibilità}:
   \begin{itemize}
   \item Recupera \texttt{inventario\_parrocchia}
   \item Calcola quantità già assegnata
   \item Verifica: \texttt{quantita} $\leq$ disponibile
   \end{itemize}

\item \textbf{Creazione Assegnazione}:
   \begin{itemize}
   \item Insert in \texttt{assegnazione\_bene}
   \item Dati: \texttt{risorsa\_id}, \texttt{beneficiario\_id}, \texttt{quantita}, \texttt{note}
   \end{itemize}

\item \textbf{Risposta}:
   \begin{itemize}
   \item Return JSON con assegnazione creata
   \end{itemize}
\end{enumerate}

\textbf{Codice rilevante}:
\begin{lstlisting}
// /api/assegnazioni/route.ts
export async function POST(request: Request) {
    // 1. Parse body
    const { risorsa_id, beneficiario_id, quantita, note } = await request.json();
    
    // 2. Validazione
    if (!risorsa_id || quantita <= 0) {
        return NextResponse.json({ error: "Parametri non validi" }, { status: 400 });
    }
    
    // 3. Autenticazione e autorizzazione
    const userResult = await getUserParrocchia(userId);
    const parrocchiaId = userResult.data.parrocchia_id;
    
    // 4. Verifica risorsa appartiene alla parrocchia
    const { data: risorsa } = await supabase
        .from("risorsa")
        .select("id, parrocchia_id")
        .eq("id", risorsa_id)
        .eq("parrocchia_id", parrocchiaId)
        .maybeSingle();
    
    // 5. Controllo disponibilita
    const { data: inventario } = await supabase
        .from("inventario_parrocchia")
        .select("quantita")
        .eq("risorsa_id", risorsa_id)
        .eq("parrocchia_id", parrocchiaId)
        .maybeSingle();
    
    const quantitaAssegnata = /* calcolo da assegnazioni esistenti */;
    const quantitaDisponibile = inventario.quantita - quantitaAssegnata;
    
    if (quantita > quantitaDisponibile) {
        return NextResponse.json(
            { error: `Quantita insufficiente. Disponibile: ${quantitaDisponibile}` },
            { status: 400 }
        );
    }
    
    // 6. Creazione assegnazione
    const { data: nuovaAssegnazione } = await supabase
        .from("assegnazione_bene")
        .insert({
            risorsa_id,
            beneficiario_id,
            quantita,
            note,
            data_assegnazione: new Date().toISOString()
        })
        .select()
        .maybeSingle();
    
    return NextResponse.json(nuovaAssegnazione, { status: 201 });
}
\end{lstlisting}

\subsubsection{Fase 5: Aggiornamento UI}

Dopo la risposta positiva dal server:
\begin{enumerate}[leftmargin=*]
\item Mostra toast di successo: "Bene assegnato!"
\item Redirect a \texttt{/beni}
\item Ricarica lista beni con quantità aggiornata
\end{enumerate}

\subsection{Validazioni e Controlli}

\subsubsection{Client-Side}
\begin{itemize}
\item Form validation (campi obbligatori)
\item Verifica quantità disponibile (basata su dati cached)
\item Feedback immediato all'utente
\end{itemize}

\subsubsection{Server-Side}
\begin{itemize}
\item Autenticazione utente
\item Autorizzazione (verifica parrocchia)
\item Validazione business logic
\item Controllo disponibilità reale (query database)
\item Transazioni atomiche
\end{itemize}

\subsection{Gestione Errori}

Esempi di errori gestiti:

\begin{lstlisting}
// 1. Quantita insufficiente
if (quantita > quantitaDisponibile) {
    return NextResponse.json(
        { error: `Quantita insufficiente. Disponibile: ${quantitaDisponibile}` },
        { status: 400 }
    );
}

// 2. Risorsa non trovata
if (!risorsa) {
    return NextResponse.json(
        { error: "Risorsa non trovata o non appartiene alla tua parrocchia." },
        { status: 404 }
    );
}

// 3. Beneficiario non trovato
if (!beneficiario) {
    return NextResponse.json(
        { error: "Beneficiario non trovato o non appartiene alla tua parrocchia." },
        { status: 404 }
    );
}
\end{lstlisting}

\section{Pattern e Convenzioni}

\subsection{Naming Conventions}

\begin{itemize}
\item \textbf{Componenti}: PascalCase (\texttt{DashboardLayout.tsx})
\item \textbf{File API}: lowercase (\texttt{route.ts})
\item \textbf{Variabili}: camelCase (\texttt{userData}, \texttt{beneficiarioId})
\item \textbf{Costanti}: UPPER\_SNAKE\_CASE (\texttt{NEXT\_PUBLIC\_SUPABASE\_URL})
\end{itemize}

\subsection{Struttura API Routes}

Ogni API route esporta funzioni HTTP standard:
\begin{itemize}
\item \texttt{GET}: Lettura dati
\item \texttt{POST}: Creazione risorse
\item \texttt{PATCH}: Aggiornamento parziale
\item \texttt{DELETE}: Eliminazione risorse
\end{itemize}

\subsection{Gestione Stato}

\begin{itemize}
\item \textbf{Client Components}: \texttt{useState} per stato locale
\item \textbf{Server Components}: Props e fetch diretti
\item \textbf{Global State}: Non utilizzato (preferiti fetch diretti)
\end{itemize}

\subsection{Error Handling}

\begin{itemize}
\item \textbf{API Routes}: Return \texttt{NextResponse.json(\{ error: "..." \}, \{ status: ... \})}
\item \textbf{Client Components}: Try-catch con toast notifications
\item \textbf{Database Errors}: Logging e messaggi user-friendly
\end{itemize}

\subsection{Type Safety}

\begin{itemize}
\item TypeScript per type checking
\item Interfacce per dati API
\item Type inference da Supabase queries
\end{itemize}

\subsection{Styling}

\begin{itemize}
\item Tailwind CSS utility classes
\item Design system con variabili CSS
\item Componenti UI riutilizzabili (shadcn/ui)
\end{itemize}

\section{Conclusioni}

L'architettura di KARIS è progettata per essere:

\begin{itemize}
\item \textbf{Scalabile}: Next.js permette crescita orizzontale
\item \textbf{Manutenibile}: Separazione chiara delle responsabilità
\item \textbf{Type-Safe}: TypeScript end-to-end
\item \textbf{User-Friendly}: Feedback immediato e UI moderna
\item \textbf{Sicura}: Autenticazione e autorizzazione a ogni livello
\end{itemize}

Il sistema utilizza tecnologie moderne e consolidate, garantendo performance, sicurezza e facilità di sviluppo.

\section{Riferimenti}

\begin{itemize}
\item Next.js Documentation: \url{https://nextjs.org/docs}
\item Supabase Documentation: \url{https://supabase.com/docs}
\item React Documentation: \url{https://react.dev}
\item Tailwind CSS Documentation: \url{https://tailwindcss.com/docs}
\end{itemize}

\end{document}

